\chapter{Technical Foundations}

This chapter fixes the graph-learning concepts used to define long-range dependency learning in later chapters.
The progression starts with graph inputs and learning tasks, then moves from local message passing to receptive-field growth.
It next states the symmetry and expressivity constraints on these models.
The final section separates information that an architecture can reach from information that a fitted model has learned to use.

\section{Graphs, Features, and Learning Tasks}

A graph is written as $G=(V,E,\lambda)$.
Here $V$ is a finite set of nodes, $E\subseteq V\times V$ is a set of directed edges, and $\lambda:V\to\mathbb{R}^{n}$ maps each node to an input feature vector.
Nodes are denoted by $u,v,w\in V$.
The feature of node $v$ is $\lambda(v)$, and $n$ is the input feature dimension.

The edge convention fixes the direction of information flow.
An edge $(u,v)\in E$ means that information can be aggregated from $u$ into $v$.
The incoming neighbourhood of $v$ is therefore
\[
N(v)=\{u\in V\mid (u,v)\in E\}.
\]
An undirected edge between $u$ and $v$ is represented by including both $(u,v)$ and $(v,u)$ in $E$.
The same notation therefore covers directed and undirected graphs without changing the aggregation convention.

Graph structure and node features provide different parts of the input.
The edge set determines which nodes are connected and the directions in which information may travel.
The feature map supplies the values attached to those nodes.
A target function may depend on either part or on their relation.

A directed path of length $d$ is a sequence $p=(v_0,\ldots,v_d)$ such that $(v_{i-1},v_i)\in E$ for every $i$ with $1\leq i\leq d$.
When a path carries information relevant to a target prediction, $v_0$ is the source node and $v_d$ is the target node.
The directed graph distance $\operatorname{dist}_G(u,v)$ is the length of a shortest directed path from $u$ to $v$, when such a path exists.
Path length and graph distance are structural properties of $G$.
Dependency length is instead a property of what the target function requires, as Chapter 3 defines in detail.

Graph learning tasks differ by where the prediction is attached.
Node-level tasks predict a value for a target node, edge-level tasks predict a value for a pair of nodes, and graph-level tasks predict one value for the whole graph.
Node classification assigns a discrete label to a target node, whereas node regression assigns a real-valued output.
Binary node classification has two labels, written as True and False in this thesis.

An inductive task requires a learned rule to generalise to graph instances or target nodes not seen during training.
Each inductive node-classification example may be represented by a graph $G$, a target node $v$, and its label.
This setting differs from transductive node classification, where training and test nodes belong to a graph whose structure is already available during training.
GLoRa, the thesis's only included paper, studies inductive binary node classification \cite{zhou2025glora}.
This scope makes the target node representation the main output used in the chapters that follow.

\section{Message-Passing Graph Neural Networks}

Message-passing graph neural networks compute node representations through repeated local updates \cite{scarselli2008graph}.
Each layer aggregates representations from the incoming neighbourhood and combines the result with the node's previous representation.
The same operation is applied at every node, which lets one set of learned parameters operate on graphs with different numbers of nodes.

Let $h_v^{(\ell)}$ denote the representation of node $v$ after message-passing layer $\ell$.
The initial representation is the input feature:
\[
h_v^{(0)}=\lambda(v).
\]
A generic layer first computes the neighbourhood message
\[
m_v^{(\ell)}=\operatorname{Agg}^{(\ell)}\left(\left\{\!\left\{h_u^{(\ell-1)}\mid u\in N(v)\right\}\!\right\}\right),
\]
where $\left\{\!\left\{\cdot\right\}\!\right\}$ denotes a multiset.
It then updates the node representation by
\[
h_v^{(\ell)}=\operatorname{Upd}^{(\ell)}\left(h_v^{(\ell-1)},m_v^{(\ell)}\right).
\]
The aggregation and update functions usually contain learnable parameters.
The layer index indicates that different layers may use different parameters.
Within a layer, the same aggregation and update functions are used at every node in $V$.

The update equation gives the node direct access to its previous representation.
An alternative convention adds self-loops and aggregates over $N(v)\cup\{v\}$.
Both conventions let a node retain or transform its own information while receiving information from its neighbours.
Under the directed-edge convention, each layer follows the available edges into the target node.

Let $L$ be the number of message-passing layers.
For node classification, a final classifier maps the target node representation to a predicted label:
\[
\hat{y}_v=\operatorname{Cls}\left(h_v^{(L)}\right).
\]
For graph-level prediction, an additional readout first forms a graph representation:
\[
h_G=\mathrm{Readout}\left(\left\{\!\left\{h_v^{(L)}\mid v\in V\right\}\!\right\}\right).
\]
The graph predictor then uses $h_G$ rather than one target node representation.
The node classifier and graph readout differ in where they attach the output, but both rely on the same sequence of local node updates.

\section{Locality and Receptive Fields}

Locality follows directly from the message-passing equations.
After one layer, $h_v^{(1)}$ can depend on $\lambda(v)$ and on the initial features of nodes in $N(v)$.
A second layer can also receive information from nodes with edges into those neighbours.
After $\ell$ layers, $h_v^{(\ell)}$ can depend only on nodes that reach $v$ by a directed path of length at most $\ell$.

This set is the receptive field of $v$ at layer $\ell$.
It is defined recursively by
\[
\mathcal{R}_0(v)=\{v\},
\]
and
\[
\mathcal{R}_{\ell}(v)=\mathcal{R}_{\ell-1}(v)\cup\{u\mid (u,w)\in E\ \textrm{for some}\ w\in\mathcal{R}_{\ell-1}(v)\}.
\]
The representation $h_v^{(\ell)}$ is a function only of the graph and features within $\mathcal{R}_{\ell}(v)$ and of the parameters in the first $\ell$ layers.
This claim follows by induction from the aggregation and update equations.

Each local layer expands the receptive field by at most one directed hop.
If $\operatorname{dist}_G(u,v)=d$ and $L<d$, then $u\notin\mathcal{R}_L(v)$.
Changing $\lambda(u)$ while holding the rest of the input fixed cannot then change $h_v^{(L)}$ or $\hat{y}_v$.
This is a structural limitation of the computation graph rather than a property of training.

When $L\geq d$, the source node $u$ lies inside the target node's receptive field.
The architecture then contains a directed computational route by which $\lambda(u)$ may affect $h_v^{(L)}$.
This condition establishes structural reachability only.
It does not establish that the aggregation and update functions preserve the relevant information or that the classifier uses it.

Depth and dependency length therefore play different roles.
Depth $L$ bounds the structural range of a local message-passing computation.
Dependency length states how far the target function requires relevant information to matter.
A model needs enough depth to reach the required information, but sufficient depth alone does not make the fitted prediction depend on that information.

\section{Symmetry and Expressivity}

Graphs have no canonical node order.
Relabelling the nodes changes their representation in memory but not the underlying graph input.
Node-level outputs should follow this relabelling, while a graph-level output should remain unchanged.

Let $\pi$ be a permutation of $V$.
The relabelled graph $G^\pi$ contains an edge $(\pi(u),\pi(v))$ whenever $(u,v)\in E$, and node $\pi(v)$ receives the feature $\lambda(v)$.
A node predictor $F$ is permutation equivariant when
\[
F(G^\pi)_{\pi(v)}=F(G)_v
\]
for every node $v$ and every permutation $\pi$.
Relabelling the input then relabels the node-level output in the same way.

A graph predictor $R$ is permutation invariant when
\[
R(G^\pi)=R(G)
\]
for every permutation $\pi$.
The graph-level output has no node index, so relabelling must leave it unchanged.

Message passing is equivariant when every node uses the same update function and the aggregation function treats its inputs as a multiset.
Changing the order in which neighbours are stored then leaves each aggregated message unchanged.
Applying a permutation-invariant readout to the final multiset of node representations produces an invariant graph representation.
These symmetry properties let the prediction depend on graph structure and features without depending on storage order.

Expressivity describes which input-output functions a model class can represent.
A target function is expressible by a model class if some setting of the model parameters realises that function.
Expressibility is therefore a property of the model class and target function, not of a particular training run.

Locality gives the first expressivity limit.
The target representation $h_v^{(L)}$ can depend only on the graph structure and features inside $\mathcal{R}_L(v)$.
An $L$-layer local message-passing model therefore cannot express a node-classification function whose output at $v$ changes when only information outside this receptive field changes.
No parameter setting can create a computational route from that information to $h_v^{(L)}$.

Even within the receptive field, a particular architecture may map inputs that the target function must distinguish to the same representation.
An aggregation function can map distinct neighbourhood multisets to the same message.
Once two inputs have produced the same intermediate representation, later layers cannot recover the distinction from that representation alone.
Permutation symmetry also prevents a model from using an arbitrary storage position as a node identifier unless that identifier is supplied as a feature.
Logical and combinatorial analyses formalise related limits on which rooted nodes and graphs message-passing architectures can distinguish \cite{barcelo2020logical}.

Learnability is a separate question.
It asks whether a training procedure can use the available data and finite resources to find parameters that realise or approximate an expressible target function.
If the target function is not expressible, no training procedure can realise it exactly.
Failure to realise it exactly therefore does not test learnability.
If the target function is expressible, successful training is still not guaranteed.
An evaluation intended to isolate learnability must therefore establish expressibility before it interprets training success or failure.

\section{From Reachability to Use}

The preceding definitions separate three claims about a target-node prediction.
Structural reachability asks whether the relevant input lies inside the target node's receptive field.
Expressibility asks whether some parameter setting in the model class can compute the required target function.
Demonstrated use is distinct from learnability: it asks whether the fitted predictor actually depends on the relevant input.

These claims are not equivalent: reachability does not establish expressibility, and expressibility does not establish demonstrated use.
A source node outside $\mathcal{R}_L(v)$ is unreachable by an $L$-layer local model.
A source node inside $\mathcal{R}_L(v)$ is reachable, but the model class may still merge inputs that the target function must distinguish.
Even when the target function is expressible, the training procedure may return parameters whose prediction ignores the source information.

Dependence concerns the behaviour of the fitted function under relevant input changes.
To test this dependence, change the relevant source or path information while holding the rest of the input fixed.
The fitted prediction should change in the way required by the target function.
Architectural depth cannot supply this evidence because depth describes possible information flow rather than the function selected by training.

An accuracy trend as graph distance or dependency length increases is likewise a behavioural symptom.
It may show where a fitted system succeeds or fails under an evaluation protocol, but it does not identify the mechanism responsible for that behaviour.
A mechanism claim requires separate evidence that the model exhibits the internal trace predicted by that mechanism.

Chapter 3 builds on this distinction by defining the path-aware dependency that an evaluation must demonstrate.
It separates graph distance and path length from the dependency length required by the target function.
The premise for that definition is the limit established here: being inside a receptive field does not demonstrate dependence.
